%%%%%%%%%%%%%%%%%%%%%%%%%%%%%%%%%%%%%%%%%%%%%%%%%%%%%%%%%%%%%%%%%%%%%%%%%%%%%%%
%% DOCTORAL EXAMINATION REPORT -- TEMPLATE
%% ---------------------------------------------------------------------------
%% Board-facing examination record for a Ph.D. by Published Works. Assembles
%% the per-paper technical examinations, the candidate's responses beneath each
%% finding, the Doctoral Examiner's dispositions, the concluding
%% recommendation, the revision schedule, declarations, and signatures.
%%
%% Compile with:  pdflatex main.tex  (run twice for the ToC and cross-refs)
%% Requirements:  a standard TeX Live installation; no external .sty files.
%%
%% HOW TO USE
%%   1. Edit the METADATA block (institution, degree, people, dates, papers).
%%   2. Set the TOGGLES (oral examination, exhibits, correspondence paper).
%%   3. For each paper, copy one of the two PAPER BLOCKS below:
%%        - Block A: paper examined through review correspondence of record;
%%        - Block B: paper examined through an itemised technical audit.
%%   4. Every finding follows the same unit:
%%        \finding{...}  statement of the finding
%%        \sev{...}      severity (omit when the finding is an endorsement)
%%        response env   the candidate's response
%%        \disposition   the Doctoral Examiner's adjudication
%%   5. Keep the Revision summary in sync with the dispositions: every
%%      "scheduled" disposition must have a line there.
%%
%% STRUCTURAL RULES ENCODED HERE
%%   - The two examination questions (Q1, Q2) are stated once in Section 1
%%     and answered explicitly, in the same words, in the recommendation.
%%   - Every finding is classified against the governing standard: evidence
%%     against doctoral competence vs. ordinary product of original research
%%     corrected through revision.
%%   - Dispositions in the revision summary use a closed vocabulary:
%%     corrected / restated / withdrawn / qualified (macros below).
%%%%%%%%%%%%%%%%%%%%%%%%%%%%%%%%%%%%%%%%%%%%%%%%%%%%%%%%%%%%%%%%%%%%%%%%%%%%%%%

\documentclass[11pt]{article}

%%%%%%%%%%%%%%%%%%%%%%%%%%%%%%%%%%% PACKAGES %%%%%%%%%%%%%%%%%%%%%%%%%%%%%%%%%%
\usepackage[utf8]{inputenc}
\usepackage[T1]{fontenc}
\IfFileExists{lmodern.sty}{\usepackage{lmodern}}{} % optional; degrades gracefully
\usepackage[margin=1in]{geometry}
\usepackage[expansion=false]{microtype} % expansion off: safe without scalable fonts
\usepackage{amsmath,amssymb}
\usepackage{enumitem}
\usepackage{xcolor}
\usepackage{fancyhdr}
\usepackage[colorlinks=true,linkcolor=black!70,urlcolor=blue!60!black,citecolor=black]{hyperref}

\setlength{\parskip}{0.45em}
\setlength{\parindent}{0em}
\setlength{\emergencystretch}{2em}

%%%%%%%%%%%%%%%%%%%%%%%%%%%%%%%%%% METADATA %%%%%%%%%%%%%%%%%%%%%%%%%%%%%%%%%%%

%% Institution
\newcommand{\institution}{Institution Name, Inc.}
\newcommand{\institutionPlace}{City, State, Country}
\newcommand{\board}{the board of \institution}

%% Degree
\newcommand{\degreeName}{Doctor of Philosophy in Mathematical Physics}
\newcommand{\degreeConcentration}{(Concentration Name)}
\newcommand{\degreeRoute}{by Published Works}

%% Governing regulations
\newcommand{\regsTitle}{\emph{Title of the Published Regulations}, YYYY Edition}
\newcommand{\regsDOI}{10.xxxx/zenodo.xxxxxxx}
\newcommand{\regsWaiverSection}{Section~X.Y} % covering-document waiver clause

%% People
\newcommand{\candidate}{A. N. Candidate}
\newcommand{\candidateShort}{A.~N.~Candidate}  % used in the running header
\newcommand{\examiner}{E. X. Aminer}
\newcommand{\examinerRole}{Doctoral Examiner}
\newcommand{\examinerTitle}{Title and Faculty, \institution}
\newcommand{\examinerAffiliation}{Secondary academic affiliation, University}

%% Dates and status
\newcommand{\examPeriod}{Month--Month YYYY}
\newcommand{\reportDate}{Month DD, YYYY}
\newcommand{\reportStatus}{Final examination report (prepared for signature and
  archival deposit)}

%% Portfolio size (words, for prose)
\newcommand{\numPapers}{five}
\newcommand{\numPapersCap}{Five}
\newcommand{\minPapers}{three}

%%%%%%%%%%%%%%%%%%%%%%%%%%%%%%%%%%% TOGGLES %%%%%%%%%%%%%%%%%%%%%%%%%%%%%%%%%%%

\newif\iforalexam       \oralexamfalse       % was an oral examination held?
\newif\ifcorrespondence \correspondencetrue  % is any paper examined through
                                             % review correspondence of record?
\newif\ifexhibits       \exhibitstrue        % include the exhibits list
\newif\ifdraftnotes     \draftnotesfalse     % show internal drafting notes

%%%%%%%%%%%%%%%%%%%%%%%%%%%%%%%%%%% MACROS %%%%%%%%%%%%%%%%%%%%%%%%%%%%%%%%%%%%

%% The unit of the examination record: finding -> severity -> response ->
%% disposition.
\newenvironment{response}{%
  \par\medskip
  \begin{list}{}{\setlength{\leftmargin}{1.6em}\setlength{\rightmargin}{0.8em}}
  \item\relax\small\textbf{Candidate's response.}\hspace{0.5em}\ignorespaces}%
  {\end{list}\par\medskip}
\newcommand{\sev}[1]{\par\smallskip{\itshape Severity: #1.}\par}
\newcommand{\disposition}[1]{\par\smallskip\noindent\textbf{\examinerRole's
  disposition.}\ \emph{#1}\par\medskip}
\newcommand{\finding}[1]{\subsubsection*{#1}}

%% Per-paper section header: \paperheader{n}{Title}{identifier}{label}
\newcommand{\paperheader}[4]{%
  \clearpage
  \section{Paper #1 examination: #2}\label{#4}
  \textbf{Manuscript.} \emph{#2.} #3.\par}

%% Closed vocabulary for the revision summary.
\newcommand{\corrected}{\emph{corrected}}
\newcommand{\restated}{\emph{restated}}
\newcommand{\withdrawn}{\emph{withdrawn}}
\newcommand{\qualified}{\emph{qualified}}
\newcommand{\rev}[2]{\item #1 --- #2.}  % \rev{what changes}{\corrected}

%% Signature block: \signatureblock{Name}{Role}{extra line(s) or empty}
\newcommand{\signatureblock}[3]{%
  \vspace{3.2em}
  \noindent\begin{tabular}{@{}p{0.55\textwidth}p{0.35\textwidth}@{}}
  \rule{0.5\textwidth}{0.4pt} & \rule{0.3\textwidth}{0.4pt}\\
  #1 & Date\\
  #2 & \\
  #3
  \end{tabular}\par}

%% Internal notes, invisible unless \draftnotestrue.
\newcommand{\note}[1]{\ifdraftnotes\textcolor{red}{\small[\,#1\,]}\fi}

%% Mathematical shorthands used across the portfolio. Extend as needed.
\newcommand{\norm}[1]{\lVert #1 \rVert}
\newcommand{\CK}{\mathrm{CK}}
\newcommand{\dist}{\operatorname{dist}}

%%%%%%%%%%%%%%%%%%%%%%%%%%%%%%%%%% PAGE STYLE %%%%%%%%%%%%%%%%%%%%%%%%%%%%%%%%%

\pagestyle{fancy}
\fancyhf{}
\fancyhead[L]{\small\itshape Doctoral Examination Report --- \candidateShort}
\fancyhead[R]{\small\thepage}
\renewcommand{\headrulewidth}{0.3pt}

\begin{document}

%%%%%%%%%%%%%%%%%%%%%%%%%%%%%%%%% TITLE PAGE %%%%%%%%%%%%%%%%%%%%%%%%%%%%%%%%%%

\begin{titlepage}
\centering
\vspace*{2em}
{\scshape\Large \institution}\\[0.25em]
{\small \institutionPlace}\\
\vspace{4em}
{\LARGE\bfseries Doctoral Examination Report}\\[1.2em]
{\large Examination for the Degree of\\ \degreeName\\ \degreeConcentration\\
 \degreeRoute}\\
\vspace{4em}
\begin{tabular}{r@{\hspace{1em}}l}
\textbf{Candidate:} & \candidate\\[0.2em]
\textbf{\examinerRole:} & \examiner\\
 & \examinerTitle\\
 & \emph{(also affiliated with \examinerAffiliation)}\\[0.2em]
\textbf{Examination period:} & \examPeriod\\[0.2em]
\textbf{Report date:} & \reportDate\\[0.2em]
\textbf{Status:} & \reportStatus\\
\end{tabular}
\vfill
\begin{minipage}{0.86\textwidth}
\small This report constitutes the official record of the doctoral examination
of the candidate's \numPapers-paper portfolio under the Institute's published
examination regulations for the degree \degreeRoute. It assembles the technical
examinations of the portfolio, the candidate's responses embedded beneath the
corresponding findings,\ifcorrespondence\ the review correspondence of record,\fi\
and the \examinerRole's concluding recommendation.
\end{minipage}
\vspace{2em}
\end{titlepage}

\tableofcontents
\clearpage

% ====================================================================
\section{The examination}
% ====================================================================

\subsection{Governance and process}\label{sec:process}

The examination was conducted under the Institute's published regulations for
the degree \degreeRoute\ (\institution, \regsTitle, DOI: \regsDOI). The process
comprises seven steps: (1) publication of the governing regulations;
(2) submission of a portfolio of published works satisfying the requirement ---
here, \numPapers\ papers against a \minPapers-preprint minimum; (3) a technical
examination of each portfolio paper; (4) the candidate's written response to
every finding; (5) independent review of the examinations, the responses, and
the underlying manuscripts by the \examinerRole, who makes the recommendation
recorded in Section~\ref{sec:recommendation}; (6) resolution by \board, from
which the candidate is recused; and (7) archival of this report and its exhibits
as a public record.
\iforalexam
An oral examination was held on \note{date}; its record is Exhibit~\note{X}.
\else
The \examinerRole\ determined that the written record provides a sufficient
basis for recommendation, and no oral examination was required.
\fi

%% Describe who prepared each technical examination and in what form.
\note{State, paper by paper, the form of the technical examination (itemised
audit, review correspondence, or other), who prepared it, and in what
capacity. Where a paper was examined through correspondence in which the
\examinerRole\ acted as reviewer, say so here and deposit the correspondence
in full as an exhibit.}

\subsection{The questions before the examiner and the governing standard}\label{sec:questions}

The examination is directed at two questions, which the concluding
recommendation of Section~\ref{sec:recommendation} answers explicitly:

\begin{enumerate}[label=\textbf{Q\arabic*.},leftmargin=3em]
\item Has the candidate demonstrated the independent ability to generate
      original, valuable doctoral-level scholarship?
\item Do the identified deficiencies materially undermine that demonstration,
      or are they appropriately addressed through revision?
\end{enumerate}

The purpose of doctoral examination is to determine whether the candidate has
demonstrated the independent ability to generate original, valuable scholarship
at the doctoral level. It is not to determine whether every result survives
subsequent criticism unchanged. Accordingly, the examination distinguishes
between deficiencies that are evidence against doctoral competence and
deficiencies that are the ordinary product of developing original research and
are appropriately corrected through revision. The object of evaluation
throughout is the candidate's demonstrated research capability; the manuscripts
are the evidence. Every finding in this report is classified on that
distinction.

\subsection{Materials examined}\label{sec:materials}

The portfolio comprises the following works, publicly deposited as preprints and
examined in the versions of record at the examination date.

\begin{enumerate}[label=\textbf{Paper \arabic*.},leftmargin=5.5em,itemsep=0.35em]
\item \emph{\note{Title of Paper 1}.} \note{Repository, DOI; earlier versions.}
\item \emph{\note{Title of Paper 2}.} \note{Repository, DOI.}
\item \emph{\note{Title of Paper 3}.} \note{Repository, DOI.}
%% add one \item per paper
\end{enumerate}

\subsection{Method and disclosure}\label{sec:method}

\note{How the technical examinations were prepared: line-by-line audit,
independent verification of every computationally checkable claim, and the
arithmetic regime used. Define ``the examiner'' and ``the \examinerRole'' as
they are used in the per-paper sections. State the standard to which the
candidate's responses were prepared and verified. State whether findings are
reproduced verbatim or in edited form, and that the manuscripts and exhibits
are incorporated by reference. If any part of the technical examination was
prepared with the assistance of software tools, including AI systems, state
which parts and how their output was checked.}

The \examinerRole\ independently reviewed the technical examinations, the
candidate's responses, and the underlying manuscripts; the terms on which he or
she adopts the findings of this report are stated in the certification
(Section~\ref{sec:signatures}).

% ====================================================================
\section{Executive summary and overall recommendation}
% ====================================================================
%% Three paragraphs: (1) the program in one paragraph, one sentence per paper;
%% (2) what the examinations found, located paper by paper; (3) the
%% recommendation in one sentence, pointing at Section \ref{sec:recommendation}.

\note{Paragraph 1 --- the research program, one sentence per paper.}

\note{Paragraph 2 --- where the examinations located the defects, paper by
paper, and how the candidate responded.}

Section~\ref{sec:recommendation} answers the two examination questions on this
record and recommends \note{outcome}, with this report and its exhibits
archived as the examination record.

% ====================================================================
\section{Portfolio assessment}\label{sec:portfolio}
% ====================================================================

\subsection{What the portfolio establishes}

\note{One paragraph per paper: the contribution that survives examination, and
what must be corrected. Close with a sentence on the coherence of the
program: whether each paper arises from an unresolved question in the
preceding work.}

\subsection{Comparison with the ordinary doctoral standard}

An ordinary Ph.D.\ is generally meant to establish that the candidate can
identify an original research problem; master the relevant literature and
mathematics; formulate definitions and conjectures; prove original results;
detect and repair errors; communicate the work coherently; and sustain
independent research. \note{State, with reference to the record, which of
these capacities are evidenced and where. Address the conditions under which
the work was produced (supervision or its absence) and whether they affect the
evidence available or the standard applied.}

\subsection{The identified deficiencies}

\note{List the principal deficiencies, paper by paper, without softening them.
Then apply the governing standard of Section~\ref{sec:questions}: are they
evidence against doctoral competence, or the ordinary product of developing
original research? Show that the valid contributions survive removal or
reconstruction of the defective portions.}

%%%%%%%%%%%%%%%%%%%%%%%%%%%%%%%%%%%%%%%%%%%%%%%%%%%%%%%%%%%%%%%%%%%%%%%%%%%%%%%
%% PAPER BLOCK A -- paper examined through review correspondence of record
%%%%%%%%%%%%%%%%%%%%%%%%%%%%%%%%%%%%%%%%%%%%%%%%%%%%%%%%%%%%%%%%%%%%%%%%%%%%%%%
\ifcorrespondence
\paperheader{1}{\note{Title of Paper 1}}{\note{Repository, DOI}}{sec:paper1}

\textbf{Form of examination.} The technical examination of Paper~1 is the
review correspondence identified in Section~\ref{sec:process}, deposited in
full as Exhibit~A and summarized here for the record. References to ``the
examiner'' in this section are to the \examinerRole\ acting as reviewer; the
dual role is disclosed, and the correspondence is deposited in full rather than
paraphrased so that the exchange can be assessed directly.

\subsection{Course of the review}

\note{Opening assessment and date.}

\emph{\note{First argument of the paper}.} \note{Objection raised; how it was
checked; what was adopted, rejected with reasons, or contributed by the
reviewer with attribution.}

\emph{\note{Second argument of the paper}.} \note{\dots}

\subsection{Assessment and disposition}

\note{What the correspondence shows about the review dynamic; how the
manuscript changed; whether any findings remain open; points recorded as
directions of continuing work rather than deficiencies.}
\fi

%%%%%%%%%%%%%%%%%%%%%%%%%%%%%%%%%%%%%%%%%%%%%%%%%%%%%%%%%%%%%%%%%%%%%%%%%%%%%%%
%% PAPER BLOCK B -- paper examined through an itemised technical audit
%% Copy this block once per paper.
%%%%%%%%%%%%%%%%%%%%%%%%%%%%%%%%%%%%%%%%%%%%%%%%%%%%%%%%%%%%%%%%%%%%%%%%%%%%%%%
\paperheader{2}{\note{Title of Paper 2}}{\note{Repository, DOI}}{sec:paper2}

\subsection*{Technical examiner's verdict}

\note{The sound core; where the defects concentrate; whether any is fatal and
to what (a theorem, a section, or the paper); the recommended architecture of
the repair.}

\begin{response}
\note{The candidate's global response: acceptance or contestation, the
verification performed, and the dependency surface into companion papers.}
\end{response}

\subsection*{Findings and responses}

\finding{Finding 1: \note{Short title}}
\note{What the manuscript claims, where. What is actually the case, with the
evidence. The consequence and the recommended repair.}
\sev{\note{fatal to X / major / major but straightforward / repairable /
minor / expositional}}
\begin{response}
\note{Accepted, contested, or accepted with refinement; what was verified;
what the revision does.}
\end{response}
\disposition{\note{Adjudication in one or two sentences. End with
``Correction scheduled.'' when a line is owed in the revision summary.}}

\finding{Finding 2: \note{Short title of an endorsement}}
\note{Endorsements carry no severity line.}
\begin{response}
\note{Response.}
\end{response}
\disposition{No correction required.}

\subsection*{Further corrections identified by the candidate}

\begin{response}
\note{Defects disclosed by the candidate beyond the examination's findings.}
\end{response}
\disposition{The self-identified corrections are added to the revision
schedule. Defects disclosed by the candidate beyond the examination's findings
are weighed as evidence under the governing standard.}

\subsection*{Corrected architecture adopted}

\note{The revised paper, as a numbered sequence of its components.}

\subsection*{Classification of results}

\emph{Sound and substantial:} \note{\dots}
\emph{Requiring correction:} \note{\dots}
\emph{Incorrect as stated:} \note{\dots}

% ====================================================================
\clearpage
\section{Concluding doctoral recommendation}\label{sec:recommendation}
% ====================================================================

\subsection{Answers to the examination questions}

\textbf{Q1. Has the candidate demonstrated the independent ability to generate
original, valuable doctoral-level scholarship?} \note{Yes / No.} \note{The
grounds, drawn from the record and not from the number of papers alone.}

\textbf{Q2. Do the identified deficiencies materially undermine that
demonstration, or are they appropriately addressed through revision?}
\note{Answer.} \note{The grounds: localisation of every fatal or major finding;
survival of the load-bearing contributions; the evidential weight of the
candidate's responses.}

\subsection{The standard applied}

The recommendation below applies the governing standard of
Section~\ref{sec:questions}: the findings of the per-paper sections were
classified by whether they are evidence against doctoral competence or the
ordinary product of developing original research, appropriately corrected
through revision, and the recommendation rests on that classification.

\subsection{Recommendation}

The \examinerRole\ recommends that the Institute:

\begin{enumerate}[label=(\arabic*),leftmargin=2.5em]
\item determine that the candidate has satisfied the requirements for
      candidacy, where procedurally distinct from conferral;
\item accept the submitted \numPapers-paper portfolio as the examined body of
      work for the degree \degreeRoute;
\item \note{optional:} waive the covering-document requirement
      (Regulations, \regsWaiverSection), \note{with grounds};
\item require correction of the defects identified in the examinations, as
      scheduled in the revision summary of Section~\ref{sec:revisions},
      incorporated into the official submitted versions of the papers;
\item document the examination and revision record, of which this report and
      its exhibits are the instrument; and
\item \note{confer the degree once the corrected portfolio satisfies the
      published requirements / alternative outcome}.
\end{enumerate}

\subsection{Final judgment}

\note{One paragraph: the portfolio against the published requirement; the
research independence demonstrated; the status of the deficiencies; what
remains before conferral.}

% ====================================================================
\section{Revision summary}\label{sec:revisions}
% ====================================================================

The following corrections, accepted in the candidate's responses, are scheduled
for incorporation into the official submitted versions of the papers before
conferral. Dispositions: \corrected\ (statement or proof repaired); \restated\
(claim reduced to what is proved); \withdrawn\ (removed, with any successor
stated as an open problem); \qualified\ (hypotheses or terminology made
explicit).

\subsection*{Paper 2 (\note{short title})}
\begin{itemize}[leftmargin=1.6em,itemsep=0.15em]
\rev{\note{What changes, in one line}}{\corrected}
\rev{\note{What changes}}{\restated}
\rev{Candidate-identified: \note{what changes}}{\corrected}
\end{itemize}

%% Repeat one \subsection* per paper with scheduled corrections.

% ====================================================================
\section{Declarations}\label{sec:declarations}
% ====================================================================
%% Standard in examination records; complete every item, writing "None" where
%% nothing applies.

\subsection*{Relationships and interests}
\note{Any prior or current professional, collaborative, institutional,
financial or personal relationship between the \examinerRole, the technical
examiner(s), the candidate, and the Institute; any dual role in the
examination; and how each was managed.}

\subsection*{Tools and assistance}
\note{Software, computational tools and AI systems used in preparing the
technical examinations or this report, the parts of the record they bear on,
and how their output was independently checked.}

% ====================================================================
\section{Certification and signatures}\label{sec:signatures}
% ====================================================================

\subsection*{\examinerRole's certification}

\note{What was independently reviewed.} By my signature I adopt the findings of
this report, the answers to the two examination questions, and the
recommendation of Section~\ref{sec:recommendation}; responsibility for the
findings as adopted rests with me; and I submit the report to \board\ for
resolution.

\signatureblock{\examiner}{\examinerRole, \institution}
  {{\small Academic affiliation: \examinerAffiliation} & \\}

\subsection*{Candidate's acknowledgment}

I acknowledge the findings of this report; I confirm that the responses
embedded herein are mine and that the corrections scheduled in
Section~\ref{sec:revisions} will be incorporated into the official submitted
versions of the papers before the portfolio is represented as final; and I
recuse myself from the board's resolution of this examination.

\signatureblock{\candidate}{Candidate}{}

% ====================================================================
\ifexhibits
\newpage
\section*{Exhibits}
\addcontentsline{toc}{section}{Exhibits}
% ====================================================================

\begin{description}[leftmargin=5.5em,style=nextline]
\ifcorrespondence
\item[Exhibit A.] Review correspondence of record for Paper~1: \note{medium,
  parties, date range, page count}.
\fi
\item[Exhibit B.] Paper 1, version of record at examination (\note{DOI}).
\item[Exhibit C.] Paper 2, version of record at examination (\note{DOI}).
%% add one \item per paper, and one per technical examination if deposited
\end{description}
\fi

\end{document}
